\section{Purpose}

\textbf{[STIPULATION]} This document treats the anomalous magnetic moments of the
leptons ($g$-2) in the FFGFT framework. The central distinction: \emph{ratios} of
$g$-2 differences are exactly derivable from the fractal structure, because all
conversion factors cancel. \emph{Absolute values} in SI units depend on the
projection factor $k_\text{geom}$ and on SI bridge constants -- they have, after
correction by $K_\text{frak}^{3/2}$, an accuracy of $0{.}01$--$0{.}15\,\%$, but
the absolute value itself is not a parameter-free geometric prediction. That is
the structurally honest status.

\section{Physical Foundations}

\textbf{[B]} The anomalous magnetic moment is the deviation from the Dirac value:
\[
  a_\ell = \frac{g-2}{2}.
\]
For a point particle Dirac predicts $g=2$, hence $a=0$. Quantum effects (vacuum
polarisation, vertex corrections) generate $a\neq0$.

\textbf{[B]} In FFGFT leptons are winding structures in the 4D torsion lattice;
the three generations carry different fractal Hausdorff exponents: electron (simple
winding, 1st generation), muon (fractally branched with exponent $p_\mu=5/3$),
tau (more branched, $p_\tau=4/3$). The anomalous moment follows from rotation of
the winding, charge distribution on it, and projection $4D\to3D$.

\section{The Geometric Formulas for Absolute Values}

\textbf{[K]} The fundamental parameters of the FFGFT $g$-2 calculation:
\[
  \varphi = \tfrac{1+\sqrt5}{2} = 1{.}618\ldots
  \quad(\text{golden ratio}),\qquad
  f = 7500\quad(\text{sub-Planck factor}),\qquad
  \xi = \tfrac43\times10^{-4}.
\]

\textbf{[K]} \textbf{Electron.} From the 3D surface of the 4D winding
$S_3=2\pi^2=19{.}739$ and the geometric projection factor
\[
  k_\text{geom} = \frac{2}{\sqrt\varphi}\times\sqrt2 = 2{.}224
\]
follows the ideal formula:
\[
  a_e^\text{(ideal)} = \frac{S_3/f}{k_\text{geom}}
  = \frac{19{.}739/7500}{2{.}224} = 1{.}184\times10^{-3}.
\]
Experiment: $a_e=1{.}1597\times10^{-3}$. Deviation ideal: $2{.}03\,\%$.

\textbf{[K]} \textbf{Muon.} The fractal additional winding:
\[
  \Delta a_\text{fractal} = \frac{4\pi}{f^{p_\mu}} = \frac{4\pi}{7500^{5/3}}
  = 4{.}373\times10^{-6},
  \qquad
  a_\mu^\text{(ideal)} = a_e + \Delta a_\text{fractal} = 1{.}188\times10^{-3}.
\]
Experiment: $a_\mu=1{.}1659\times10^{-3}$. Deviation ideal: $1{.}89\,\%$.

\section{Resolution of the Systematic Deviation: \texorpdfstring{$K_\text{frak}^{3/2}$}{K\_frak\textasciicircum 3/2}}

\textbf{[K]} The $\sim2\,\%$ deviation of the ideal formulas is not an open
question but is explained by the fractal correction. The effective projection
factor:
\[
  k_\text{eff}
  = \frac{k_\text{geom}}{K_\text{frak}^{3/2}}
  = \frac{2{.}224}{(1-100\xi)^{3/2}}
  = \frac{2{.}224}{0{.}98007}
  = 2{.}2692.
\]
The exponent $3/2=D/2$ describes how the fractal structure acts on the projection
across three spatial dimensions. $K_\text{frak}=1-100\xi$ is the same correction
as in A040 -- no new parameter.

\textbf{[K]} With the corrected factor:

\begin{center}
\renewcommand{\arraystretch}{1.3}
{\small
\begin{tabular}{lccccc}
\toprule
Lepton & T0 ideal & T0 corrected & Experiment & Dev.\ ideal & Dev.\ corr. \\
\midrule
Electron & $1{.}184\times10^{-3}$ & $1{.}1598\times10^{-3}$ & $1{.}1597\times10^{-3}$ & 2{.}03\,\% & \textbf{0{.}01\,\%} \\
Muon     & $1{.}188\times10^{-3}$ & $1{.}1642\times10^{-3}$ & $1{.}1659\times10^{-3}$ & 1{.}89\,\% & \textbf{0{.}15\,\%} \\
Tau      & $1{.}269\times10^{-3}$ & $1{.}2454\times10^{-3}$ & (not measured)          & --         & Prediction \\
\bottomrule
\end{tabular}
}
\renewcommand{\arraystretch}{1.0}
\end{center}

\textbf{[B]} \textbf{Methodological clarification.} The corrected absolute value
is a consequence of the FFGFT-internally consistent $K_\text{frak}$ -- not a fit.
The correction agrees with the experimentally reconstructed value
$k_\text{rec}=S_3/(f\cdot a_e^\text{exp})=2{.}2696$ to $0{.}014\,\%$.
But: $k_\text{geom}$ and $f$ depend on SI bridge constants and measured quantities.
An absolute SI value without any measured anchor -- as the SM can achieve with
$\alpha$ -- is not available at the same level in FFGFT (the same holds for
masses, A130, A080).

\section{Ratios: Parameter-Free and Exact}

\textbf{[B]} In contrast to absolute values, in ratios of $g$-2 differences all
conversion factors cancel:
\[
  \frac{\Delta a(\tau-\mu)}{\Delta a(\mu-e)}
  = \frac{4\pi/f^{4/3}-4\pi/f^{5/3}}{4\pi/f^{5/3}}
  = \frac{f^{5/3}}{f^{4/3}} - 1
  = f^{1/3} - 1.
\]
Numerically: $f^{1/3}=7500^{1/3}=19{.}57$, hence
\[
  \boxed{\;\frac{\Delta a(\tau-\mu)}{\Delta a(\mu-e)} = f^{1/3}-1 = 18{.}57\;}
\]
No $k_\text{geom}$, no $K_\text{frak}$, no SI factor, no $\alpha$. The ratio
depends exclusively on $f=7500$ -- and $f$ is the sub-Planck factor $30000/4$,
following directly from $\xi=4/30000$.

\textbf{[S]} \textbf{Structural parallel to the Koide formula.} The Koide formula
$(m_e+m_\mu+m_\tau)/(\sqrt{m_e}+\sqrt{m_\mu}+\sqrt{m_\tau})^2\approx2/3$
is confirmed to $0{.}001\,\%$ because there ratios rather than absolute values are
compared. The same structure operates here: the exponent step $\Delta p=1/3$ of
the masses (A100) appears directly as $q_\mu=5/3$, $q_\tau=4/3$, $\Delta q=1/3$
in the $g$-2 exponents. No coincidence -- the same geometric origin.

\section{The Bridge Formula: Masses and \texorpdfstring{$g$-2}{g-2} Connected}

\textbf{[B]} From the mass formulas: $m_\tau/m_\mu=(125/144)\cdot f^{1/3}$.
From the $g$-2 formulas: $\Delta a(\tau-e)/\Delta a(\mu-e)=f^{1/3}$.
Combining gives the bridge formula:
\[
  \boxed{\;\frac{\Delta a(\tau-e)}{\Delta a(\mu-e)} = \frac{144}{125}\cdot\frac{m_\tau}{m_\mu}\;}
\]
The parameter $f$ cancels. No $\alpha$, no $k_\text{eff}$, no $K_\text{frak}$.
The prediction follows solely from the shared $1/3$ step size of masses and $g$-2.

\textbf{[K]} \textbf{Tau prediction from the bridge formula.} With measured values
$a_e^\text{exp}=1{.}160\times10^{-3}$,
$\Delta a(\mu-e)^\text{exp}=6{.}269\times10^{-6}$,
$m_\tau/m_\mu=16{.}817$:
\begin{align*}
  \Delta a(\tau-e) &= \frac{144}{125}\cdot16{.}817\cdot6{.}269\times10^{-6}
  = 1{.}214\times10^{-4},\\
  a_\tau^\text{(bridge)} &= a_e + \Delta a(\tau-e)
  = \boxed{1{.}282\times10^{-3}}.
\end{align*}
The SM calculates $a_\tau^\text{SM}=1{.}177\times10^{-3}$. Difference: $9\,\%$.
Testable at Belle II ($\sim2027$).

\section{Fair Comparison: SM and FFGFT}

\textbf{[S]} \textbf{The SM circularity.} The SM extracts $\alpha$ from the
$a_e$ measurement and then computes $a_e$ back. For electron and muon this is a
circular fit -- a precision calculation, not an independent prediction. FFGFT
derives $\alpha$ from $\xi$ alone: $\alpha=\xi\cdot E_0^2$ with
$E_0=7{.}398\,\text{MeV}$ gives $1/\alpha_\text{T0}=137{.}035$ -- $0{.}0005\,\%$
from experiment (A130). That is a genuine, independent prediction, though at a
different precision level than the SM perturbation calculation.

\textbf{[S]} \textbf{Where the theories diverge:} Not at $e$ or $\mu$ (both
theories lie close to experiment there), but at the \emph{tau}. T0:
$a_\tau\approx1{.}282\times10^{-3}$, SM: $1{.}177\times10^{-3}$ -- factor 7
larger deviation than the muon discrepancy. That is the decisive test.

\section{Verification Script}

\begin{itemize}[leftmargin=2em,itemsep=2pt,topsep=2pt]
  \item Verification of all formulas: ideal and corrected absolute values, ratio
    $f^{1/3}-1$, bridge formula with experimental values, tau prediction:
    \texttt{python/A\_Serie\_Skripte/a138\_g2.py}
\end{itemize}

\section{Symbol Glossary}

Symbols used throughout the entire A-series are explained in A010.
Specific to this document:

\begin{center}
\renewcommand{\arraystretch}{1.3}
{\small\begin{tabular}{lp{9cm}}
\toprule
Symbol & Meaning \\
\midrule
$a_\ell = (g-2)/2$ & Anomalous magnetic moment of lepton $\ell$ \\
$f = 30000/4 = 7500$ & Sub-Planck factor; follows directly from $\xi=4/30000$ \\
$k_\text{geom} = (2/\sqrt\varphi)\sqrt2$ & Geometric projection factor \\
$k_\text{eff} = k_\text{geom}/K_\text{frak}^{3/2}$ & Effective projection factor
  after fractal correction \\
$S_3 = 2\pi^2$ & Surface of the 3-sphere (volume of the 4D unit ball) \\
$p_\mu = 5/3$, $p_\tau = 4/3$ & Fractal Hausdorff exponents of the
  generation structure \\
\bottomrule
\end{tabular}}
\renewcommand{\arraystretch}{1.0}
\end{center}

\section{Open Questions}

\textbf{The absolute SI values depend on SI bridge constants and measured
quantities.} The correction $K_\text{frak}^{3/2}$ explains the systematic
deviation; but the geometric projection factor $k_\text{geom}$ is not derived
at the same level as the pure ratios.

\textbf{The tau prediction ($a_\tau\approx1{.}282\times10^{-3}$) has not yet
been measured.} Belle II will test it. Until then it is a prognosis --
precisely formulated, but not a result.

\textbf{The exponent $3/2$ for the fractal correction is geometrically motivated
($D/2=3/2$) but not uniquely proved.} That exactly $D/2$ and not some other
exponent appears is a consistency observation, not a derivation.

\section{Supersedes}

This document subsumes: Dok.~018 (all three versions 018-10/11/12, anomalous
magnetic moments of the leptons: geometric formulas, $K_\text{frak}^{3/2}$
correction, ratio predictions, bridge formula masses $\leftrightarrow$ $g$-2,
tau prediction, fair SM comparison). Not taken over: narrative passages and
cosmological generalisations from Dok.~018.

\section{Use of AI Tools}

This document was created with the assistance of AI tools. The questions,
stipulations, and all substantive decisions come from the author; the tools
formulate, compute, and create verification scripts.
Responsibility for content and errors rests entirely with the author.
Full wording of the rules in A010.
